\clearpage
\pdfbookmark[0]{Vorwort.}{praefatio}
\begin{center}\Large VORWORT.\end{center}
\medskip
\SourceLeaf{010}Nachdem \textsc{Kepler}s Genie die Gesetze der Planetenbewegung entdeckt hatte, standen genügend Hilfsmittel zur Verfügung, um die Elemente der einzelnen Planeten aus den Beobachtungen zu gewinnen. \textsc{Tycho Brahe}, der die praktische Astronomie auf eine bis dahin unbekannte Höhe geführt hatte, hatte sämtliche Planeten über eine lange Reihe von Jahren mit solcher Sorgfalt und Ausdauer beobachtet, dass \textsc{Kepler}, dem würdigsten Erben eines solchen Schatzes, nur noch die Auswahl der Beobachtungen blieb, die für den jeweils gesteckten Zweck geeignet schienen. Wesentlich erleichtert wurde diese Arbeit zudem durch die mittleren Bewegungen der Planeten, die schon lange zuvor aus den ältesten Beobachtungen mit größter Genauigkeit bestimmt worden waren.

Die Astronomen, die nach \textsc{Kepler} versuchten, die Planetenbahnen mithilfe neuerer oder besserer Beobachtungen noch genauer zu vermessen, konnten auf dieselben oder sogar auf noch mehr Hilfsmittel zurückgreifen. Die Aufgabe bestand nicht mehr darin, völlig unbekannte Elemente zu ermitteln. Die bekannten Elemente mussten nur geringfügig berichtigt und enger eingegrenzt werden.

Das vom großen \textsc{Newton} entdeckte Prinzip der allgemeinen Gravitation eröffnete ein ganz neues Gebiet. Es zeigte, dass sich \emph{alle} Himmelskörper, deren Bewegung allein von der Kraft der Sonne bestimmt wird, notwendig
\SourceLeaf{011}denselben Gesetzen fügen müssen, die nach \textsc{Kepler}s Befund für die fünf Planeten galten, sofern man diese Gesetze nur geringfügig abändert. Gestützt auf die Beobachtungen hatte \textsc{Kepler} erklärt, dass die Bahn jedes Planeten eine Ellipse ist, auf der um die Sonne, die einen der Brennpunkte einnimmt, gleichmäßig Flächen überstrichen werden, und zwar so, dass sich die Umlaufzeiten auf verschiedenen Ellipsen wie die \(3/2\)-Potenzen der großen Halbachsen verhalten. Ausgehend vom Prinzip der allgemeinen Gravitation wies \textsc{Newton} a priori nach, dass sich alle Körper, die der Anziehungskraft der Sonne unterliegen, auf Kegelschnitten bewegen müssen. Die Planeten zeigen uns eine Art von Kegelschnitten, nämlich Ellipsen. Die übrigen Arten, Parabeln und Hyperbeln, müssen als ebenso möglich gelten, sofern sich Körper mit der jeweils erforderlichen Geschwindigkeit unter dem Einfluss der Sonnenkraft bewegen. Ferner wies er nach, dass die Sonne stets einen Brennpunkt des Kegelschnitts einnimmt, dass die Flächen, die ein und derselbe Körper in verschiedenen Zeitspannen um die Sonne überstreicht, zu diesen Zeitspannen proportional sind und dass sich schließlich die Flächen, die verschiedene Körper in gleichen Zeitspannen um die Sonne überstreichen, wie die Quadratwurzeln der Halbparameter ihrer Bahnen verhalten. Das letzte dieser Gesetze stimmt bei elliptischer Bewegung mit \textsc{Kepler}s letztem Gesetz überein und gilt auch für parabolische und hyperbolische Bewegung. Auf diese lässt sich Keplers Gesetz mangels Umläufen nicht anwenden. Damit war ein Faden gefunden, an dem man sich in das bisher unzugängliche Labyrinth der Kometenbewegungen wagen konnte. Das gelang so gut, dass eine einzige Hypothese, nämlich die Annahme parabolischer Bahnen, zur Erklärung der Bewegungen aller Kometen genügte, die genau beobachtet worden waren. Damit hatte das System der allgemeinen Gravitation der Analysis neue und höchst glänzende Triumphe verschafft. Die Kometen, die bis dahin immer ungebändigt geblieben waren oder sich bald wieder aufrührerisch und widerspenstig gezeigt hatten, wenn sie besiegt schienen, ließen sich die Zügel anlegen. Aus Feinden zu Gästen geworden, setzten sie ihren Weg auf den von der Rechnung vorgezeichneten Bahnen fort und fügten sich getreulich denselben ewigen Gesetzen wie die Planeten.

Schon die Bestimmung parabolischer Kometenbahnen aus Beobachtungen brachte weit größere Schwierigkeiten mit sich als die Bestimmung elliptischer Planetenbahnen. Das lag vor allem daran, dass Kometen nur während einer kürzeren
\SourceLeaf{012}Zeitspanne sichtbar waren. Der Mathematiker konnte also keine Beobachtungen auswählen, die sich für den einen oder anderen Zweck besonders eigneten. Er war auf diejenigen angewiesen, die der Zufall ihm geboten hatte, sodass die bei den Planetenrechnungen eingesetzten besonderen Verfahren kaum je angewandt werden konnten. Der große \textsc{Newton} selbst, der führende Mathematiker seines Jahrhunderts, machte aus der Schwierigkeit des Problems kein Geheimnis, ging aber, wie zu erwarten war, auch aus diesem Wettstreit als Sieger hervor. Viele Mathematiker nach \textsc{Newton} widmeten sich demselben Problem, freilich mit unterschiedlichem Erfolg, doch bleibt heute kaum noch etwas zu wünschen übrig.

Dabei ist zu bedenken, dass auch bei diesem Problem die Kenntnis eines Elements des Kegelschnitts die Schwierigkeit auf sehr willkommene Weise verringert. Schon mit der Annahme einer parabolischen Bahn wird die große Achse als unendlich groß angesetzt. Denn alle Parabeln unterscheiden sich, wenn man von ihrer Lage absieht, allein durch den größeren oder kleineren Abstand des Scheitels vom Brennpunkt, während Kegelschnitte im Allgemeinen eine unendlich größere Vielfalt zulassen. Es gab freilich keinen hinreichenden Grund für die Annahme, dass die Kometenbahnen vollkommen parabolisch seien. Vielmehr muss es als unendlich unwahrscheinlich gelten, dass die Natur einer solchen Annahme je entsprochen hat. Es stand jedoch fest, dass die Erscheinungen eines Himmelskörpers auf einer Ellipse oder Hyperbel, deren große Achse im Verhältnis zum Parameter sehr groß ist, in der Nähe des Perihels nur sehr wenig von der Bewegung auf einer Parabel mit demselben Abstand des Scheitels vom Brennpunkt abweichen. Die Abweichung wird umso geringer, je größer jenes Verhältnis der Achse zum Parameter ist. Ferner hatte die Erfahrung gezeigt, dass zwischen der beobachteten und der für eine parabolische Bahn berechneten Bewegung kaum je größere Unterschiede übrig blieben, als sich den hier meist recht erheblichen Beobachtungsfehlern selbst mit Sicherheit zuschreiben ließen. Die Astronomen hielten es daher für angebracht, bei der Parabel zu bleiben. Mit Recht, denn es fehlte jedes Hilfsmittel, mit dem man hinreichend sicher hätte ermitteln können, ob und in welchem Maß Abweichungen von der Parabel vorlagen. Eine Ausnahme bildet der berühmte \textsc{Halley}sche Komet, der eine sehr lang gestreckte Ellipse beschreibt. Weil man ihn wiederholt beobachtet hat, wenn er
\SourceLeaf{013}zum Perihel zurückkehrte, ist uns seine Umlaufzeit bekannt geworden. Ist daraus die große Achse bekannt, so dürfte die Berechnung der übrigen Elemente kaum schwieriger sein als die Bestimmung einer parabolischen Bahn. Zu erwähnen ist auch, dass Astronomen bei einigen anderen, etwas länger beobachteten Kometen versucht haben, die Abweichung von der Parabel zu bestimmen. Alle zu diesem Zweck vorgeschlagenen oder angewandten Verfahren beruhen jedoch auf der Annahme, dass die Abweichung von der Parabel gering ist. Bei diesen Versuchen lieferte daher die zuvor berechnete Parabel bereits Näherungswerte für die einzelnen Elemente, abgesehen von der großen Achse oder der von ihr abhängigen Umlaufzeit. Diese Werte bedurften nur geringer Berichtigungen. Zudem muss man einräumen, dass alle diese Versuche kaum je eine sichere Entscheidung ermöglichten, wenn man vielleicht den Kometen von 1770 ausnimmt.

Sobald feststand, dass sich die Bewegung des 1781 entdeckten neuen Planeten mit der parabolischen Hypothese nicht vereinbaren ließ, begannen die Astronomen, ihm eine Kreisbahn anzupassen. Das lässt sich mit einer sehr leichten und einfachen Rechnung bewerkstelligen. Glücklicherweise war die Bahn dieses Planeten nur mäßig exzentrisch. So lieferten die ·durch jene Annahme gewonnenen Elemente wenigstens eine grobe Näherung, auf der sich dann die Bestimmung der elliptischen Elemente aufbauen ließ. Mehrere andere sehr günstige Umstände kamen hinzu. Die langsame Bewegung des Planeten und die sehr geringe Neigung seiner Bahn gegen die Ebene der Ekliptik vereinfachten die Rechnungen erheblich und gestatteten die Anwendung besonderer Verfahren, die für andere Fälle ungeeignet waren. Zugleich zerstreuten sie die Sorge, dass der Planet, in den Strahlen der Sonne verschwunden, sich später der Suche entziehen könnte. Unter anderen Umständen, besonders wenn sein Licht außerdem schwächer gewesen wäre, hätte diese Sorge die Gemüter durchaus beunruhigen können. Eine genauere Bestimmung der Bahn ließ sich somit gefahrlos aufschieben, bis man aus zahlreicheren und zeitlich weiter auseinanderliegenden Beobachtungen diejenigen auswählen konnte, die dem Zweck am besten zu entsprechen schienen.

\SourceLeaf{014}In allen Fällen also, in denen die Bahnen von Himmelskörpern aus Beobachtungen abzuleiten waren, gab es beachtliche Vorteile, die die Anwendung besonderer Verfahren nahelegten oder wenigstens ermöglichten. Der wichtigste davon war, dass man durch hypothetische Annahmen schon Näherungswerte für einige Elemente gewinnen konnte, ehe man sich an die Berechnung der elliptischen Elemente machte. Dennoch erscheint es recht erstaunlich, dass die allgemeine Aufgabe

\begin{quote}\itshape
Die Bahn eines Himmelskörpers ohne jede hypothetische Annahme aus Beobachtungen zu bestimmen, die nur einen kurzen Zeitraum umfassen und daher keine für die Anwendung besonderer Verfahren geeignete Auswahl zulassen
\end{quote}
\noindent bis zum Beginn dieses Jahrhunderts beinahe ganz vernachlässigt oder zumindest von niemandem ernsthaft und angemessen behandelt worden ist. Ihre Schwierigkeit und Eleganz hätten sie gewiss den Theoretikern empfehlen können, auch wenn ihr großer Nutzen unter den Praktikern noch nicht feststand. Es hatte sich nämlich allgemein die Meinung durchgesetzt, dass eine solche vollständige Bestimmung aus Beobachtungen, die auf eine kürzere Zeitspanne beschränkt sind, \emph{unmöglich} sei. Diese Meinung war schlecht begründet, da inzwischen mit größter Sicherheit erwiesen ist, dass sich die Bahn eines Himmelskörpers ohne jede hypothetische Annahme schon aus guten Beobachtungen weniger Tage hinreichend genau bestimmen lässt.

Im September 1801, als ich mit einer ganz anderen Arbeit beschäftigt war, kam ich auf einige Gedanken, die zur Lösung der großen Aufgabe geeignet schienen, von der ich gesprochen habe. In einem solchen Fall lassen wir, um uns nicht zu sehr von einer fesselnden Untersuchung ablenken zu lassen, nicht selten Gedankenverbindungen unbeachtet verloren gehen, die bei näherer Prüfung reiche Früchte hätten tragen können. Vielleicht hätte auch diesen kleinen Einfällen dasselbe Schicksal gedroht, wären sie mir nicht gerade zu einer Zeit gekommen, wie man sie sich günstiger für ihre Bewahrung und Weiterentwicklung nicht hätte wünschen können. Um dieselbe Zeit war nämlich die Nachricht in aller Munde, dass am 1.~Januar jenes Jahres an der Sternwarte von Palermo ein neuer Planet entdeckt worden war. Bald wurden auch die Beobachtungen öffentlich bekannt, die der hervorragende Astronom \textsc{Piazzi} von
\SourceLeaf{015}jenem Tag bis zum 11.~Februar angestellt hatte. Damals bot sich in jener schwierigen Lage und unter solch dringendem Handlungsdruck die Gelegenheit, die Bedeutung dieser Aufgabe so deutlich wie nur möglich vor Augen zu führen. In den Annalen der Astronomie findet sich gewiss keine gleich gewichtige Gelegenheit, und eine gewichtigere ließe sich kaum denken. Denn alle Hoffnung, jenes planetarische Atom nach fast einem Jahr unter zahllosen kleinen Sternen am Himmel wiederzufinden, hing einzig von einer hinreichend genauen Kenntnis seiner Bahn ab, die allein auf diesen wenigen Beobachtungen beruhen konnte. Hätte ich je zu einem günstigeren Zeitpunkt erproben können, was meine kleinen Einfälle in der Praxis taugten, als indem ich sie damals zur Bestimmung der Bahn der Ceres einsetzte? Der Planet hatte in jenen 41 Tagen geozentrisch nur einen Bogen von drei Grad beschrieben und musste nach Ablauf eines Jahres in einer sehr weit davon entfernten Himmelsgegend gesucht werden. Diese erste Anwendung des Verfahrens erfolgte im Oktober 1801. Gleich die erste klare Nacht, in der man den Planeten anhand der so gewonnenen Zahlen suchte\footnote{Am 7.~Dezember 1801 durch den angesehenen von \textsc{Zach}.}, führte den Flüchtling der Beobachtung wieder zu. Drei weitere seit jener Zeit entdeckte Planeten boten neue Gelegenheiten, die Wirksamkeit und Allgemeinheit des Verfahrens zu prüfen und zu bestätigen.

Unmittelbar nach der Wiederauffindung der Ceres wünschten mehrere Astronomen, dass ich die bei diesen Rechnungen angewandten Verfahren veröffentlichte. Mehreres hinderte mich damals daran, diesen freundlichen Bitten nachzukommen. Ich hatte andere Verpflichtungen, wollte die Sache eines Tages ausführlicher behandeln und erwartete vor allem, dass die weitere Arbeit an dieser Untersuchung die verschiedenen Teile der Lösung zu größerer Allgemeinheit, Einfachheit und Eleganz führen würde. Da mich diese Hoffnung nicht getäuscht hat, glaube ich keinen Grund zu haben, den Aufschub zu bereuen. Die anfangs benutzten Verfahren wurden immer wieder so vielfach und so tiefgreifend verändert, dass zwischen der damaligen Berechnung der Bahn der Ceres und dem Verfahren, das
\SourceLeaf{016}in diesem Werk dargelegt wird, kaum eine Spur von Ähnlichkeit übrig geblieben ist. Eine vollständige Darstellung aller dieser nach und nach immer weiter vervollkommneten Untersuchungen lag zwar außerhalb meiner Absicht, doch hielt ich es bei mehreren Gelegenheiten, besonders wenn es um ein gewichtigeres Problem ging, für angebracht, auch die früheren Verfahren wenigstens teilweise darzustellen. Über die Lösungen der Hauptprobleme hinaus habe ich in diesem Werk vieles behandelt, was mir bei meiner recht langen Beschäftigung mit den Bewegungen der Himmelskörper auf Kegelschnitten wegen seiner analytischen Eleganz oder vor allem wegen seines praktischen Nutzens besondere Aufmerksamkeit zu verdienen schien. Dabei habe ich stets den Gegenständen oder Verfahren, die mir eigen sind, mehr Sorgfalt gewidmet. Bekanntes habe ich nur kurz berührt, soweit es der Zusammenhang der Sache zu erfordern schien.

Das ganze Werk gliedert sich demnach in zwei Teile. Im ersten Buch werden die Beziehungen zwischen den Größen entwickelt, von denen die Bewegung der Himmelskörper um die Sonne nach den Gesetzen \textsc{Kepler}s abhängt. In den beiden ersten Abschnitten geht es um die Beziehungen für einen einzelnen, für sich betrachteten Ort. Im dritten und vierten Abschnitt geht es um die Beziehungen, bei denen mehrere Orte miteinander verglichen werden. Die ersten beiden Abschnitte stellen die gebräuchlichen Verfahren und vor allem weitere Verfahren dar, die meines Erachtens für den praktischen Gebrauch bei Weitem vorzuziehen sind und von den bekannten Elementen zu den Erscheinungen führen. Die beiden folgenden Abschnitte behandeln viele sehr wichtige Aufgaben, die den Weg zu den umgekehrten Operationen bahnen. Die Erscheinungen selbst entstehen aus einer kunstvollen und verwickelten Verflechtung der Elemente. Man muss diese Art der Verflechtung genauer durchschaut haben, ehe man mit Aussicht auf Erfolg die Fäden entwirren und das Werk in seine Elemente auflösen kann. Im ersten Buch werden also die Werkzeuge und Hilfsmittel bereitgestellt, mit denen dann im zweiten Buch diese schwierige Aufgabe selbst bewältigt wird. Der größte Teil der Arbeit besteht dann darin, jene Hilfsmittel richtig zusammenzutragen, in eine geeignete Ordnung zu bringen und auf das gesteckte Ziel auszurichten.

Die meisten der wichtigeren Aufgaben wurden durch geeignete
\SourceLeaf{017}Beispiele veranschaulicht, die, wo immer es möglich war, tatsächlichen Beobachtungen entnommen sind. So wird das Vertrauen in die Wirksamkeit der Verfahren gestärkt und ihre Anwendung deutlicher vor Augen geführt. Ich hoffe zudem, dass auch weniger Geübte vom Studium dieser Gegenstände nicht abgeschreckt werden, die ohne Zweifel den fruchtbarsten und schönsten Teil der theoretischen Astronomie ausmachen.

Ich schrieb dies in Göttingen am 28.~März 1809.
